\section{Related Work}

Three lines of work lead to our question. Visual explanations ask which image content is associated with an answer. Causal tracing tests whether changing an internal state changes that answer. Sparse circuit methods make selected parts of the computation easier to inspect. Their observations are useful, but an associated region, a consequential state, and a faithful account of the score change caused by masking the evidence are different findings.

\stitle{Visual relevance and regional evidence.}
VQA benchmarks report whether an answer is correct, but language priors can support the same answer without the relevant pixels \citep{antol2015vqa,goyal2017vqav2,agrawal2018vqacp}. Saliency and attention methods identify image content associated with a prediction, while human-supervised explanations and counterfactual images test alignment or output sensitivity \citep{selvaraju2017gradcam,jain2019attention,selvaraju2019hint,goyal2019counterfactual}. These are useful observations, yet a plausible map need not reflect the learned decision function \citep{adebayo2018sanity}, and masking any region can damage the image. For our question, the relevant input-side comparison is therefore not simply clean versus masked: it is the answer-score damage from masking annotated evidence versus masking comparable non-evidence content on the same image. Once this additional damage is identified, the next question is where the score lost under the evidence mask can be recovered inside the model.

\stitle{Internal causal localization.}
Causal tracing and activation patching restore clean internal states during a corrupted computation and measure how much of the lost answer support returns \citep{vig2020causalMediation,meng2022rome,wang2023ioi}. Their conclusions depend on the chosen corruption and answer metric \citep{zhang2024patching}. In VLMs, NOTICE uses paired semantic images and text replacements to identify consequential attention heads \citep{golovanevsky2025notice}; Fine-grained Cross-modal Causal Tracing tests visual and text positions and several component types across layers to study object representations \citep{li2026causalTracing}. These studies establish that changing selected multimodal states can alter an answer. We build on this intervention logic by fixing a human-localized visual effect against a same-area control and testing both clean-state restoration and corrupted-state replacement. This paired design also creates a measured effect for evaluating a proposed sparse or compact description.

\stitle{Sparse descriptions and causal fidelity.}
Sparse autoencoders describe dense activations with a small number of learned features; transcoders instead use sparse features to approximate a model's feed-forward update \citep{cunningham2024sparseAutoencoders,dunefsky2024transcoders}. A per-layer transcoder approximates one layer's update, whereas a cross-layer transcoder can describe contributions to later layers \citep{dunefsky2024transcoders,ameisen2025circuittracing}. These are separately fitted measurement tools, not components that the base VLM uses to answer. VLM studies have traced visual concepts with attribution graphs and tested feature sets through steering or circuit patching \citep{yang2026vlmcircuittracing,zhou2026sparsevisual}. Circuit-faithfulness research already distinguishes structural overlap from functional reproduction of a circuit's effect \citep{hanna2024faithfulness}. We bring that criterion to a human-localized VQA effect: can selected features reproduce the answer-score change associated with removing the annotated visual evidence? Native circuits can use different components for analogous tasks across modalities \citep{sameTaskDifferentCircuits}; low-dimensional subspace interventions test distributed variables but require care when interpreted as natural mechanisms \citep{geiger2024das,makelov2024subspace}. We compare these different candidate descriptions with a complete reference at each tested interface rather than align feature identities across models. The next section introduces the base-model states and fitted interfaces; the Method defines how their candidate changes are tested.
